\paragraph{\trFloerCohomology。}\label{sub:floer-cohomology}
取定某个特征为 2 的域 $\KK$。对 $M$ 内的每一对\trExact\trLagrange\trSubmanifold{}$(L_0,L_1)$，
都有一个\trFloerCohomology{}群 $HF^{\mathrm{pr}}(L_0,L_1)$，它是一个有限维 $\KK$-\trVectorSpace
（pr 意即 “preliminary”（初步），用以区别于后文将要引入的更精细的理论——在后者中%
$\mathrm{char}(\KK)$ 可以任意，而且\trFloer{}群是有\trGraded{}的）。对任何带\trLagrange\trLabel{}的%
\trPointedBoundary\trRiemann\trSurface{}$S$，都可以联系上一个\trHomomorphism
\begin{equation}\label{eq:tqft-operation}
\Phi_S\colon
\bigotimes_{\zeta^+\in\Sigma^+}HF^{\mathrm{pr}}(L_{\zeta^+,0},L_{\zeta^+,1})
\longrightarrow
\bigotimes_{\zeta^-\in\Sigma^-}HF^{\mathrm{pr}}(L_{\zeta^-,0},L_{\zeta^-,1})
\end{equation}
（\trContravariant{}约定可能显得奇怪，但请记得我们处理的是\trFloerCohomology）。这些映射遵从%
\trOpenString{} TQFT 公理的一个版本（例如见 [122]），其中最重要的是如下%
\textbf{\trCompositionLaw}（composition law）。设 $S_1$、$S_2$ 是两个带\trLagrange\trLabel{}的%
\trPointedBoundary\trRiemann\trSurface。假设存在无穷远点 $\zeta_1^+\in\Sigma_1^+$、
$\zeta_2^-\in\Sigma_2^-$，它们配有同一（有序）\trLagrange\trSubmanifold{}对。把它们如%
\cref{fig:gluing-surfaces} 那样\trGluing{}起来，就得到另一个\trPointedBoundary\trRiemann\trSurface%
$S=S_1\#S_2$（由于 TQFT 的另一条性质——\emph{\trTopology\trInvariance}——具体怎样\trGluing{}%
并不要紧）。$S$ 带有诱导的\trLagrange\trLabel，并且
\begin{equation}\label{eq:composition-law}
\Phi_S=(\mathrm{Id}\otimes\cdots\otimes\Phi_{S_1}\otimes\cdots\otimes\mathrm{Id})\circ
(\mathrm{Id}\otimes\cdots\otimes\Phi_{S_2}\otimes\cdots\otimes\mathrm{Id}),
\end{equation}
其中的规则是：把 $\Phi_{S_2}$ 的输出中的 $\zeta_2^-$-因子插入 $\Phi_{S_1}$ 的
$\zeta_1^+$-输入分量。
